\subsection{The complete \texorpdfstring{$C_6\times C_{12}$}{C6 x C12} case}
\begin{proposition}\label{prop:decic612-complete}
The noncuspidal closed points of degree five on
$X_1(6,12)_{\Q(\zeta_6)}$ form one orbit under the pairing-preserving
level automorphisms. They give the curve of Proposition~\ref{prop:decic612}.
In particular every decic elliptic curve containing
$C_6\times C_{12}$ has full torsion $C_6\times C_{12}$.
\end{proposition}
\begin{proof}
Put $F=\Q(z)$, $z^2-z+1=0$, $X=X_1(6,12)$ and
$J=\operatorname{Jac}(X)$. We have $g(X)=9$, $J(F)$ has rank zero,
and $\operatorname{gon}_F X\ge6$ by \cite{DS,NV}.
There are no noncuspidal $F$-points: at $7$ there are no
noncuspidal rational residue points, and
Lemma~\ref{lem:decic-cuspidal-reduction} excludes cuspidal reduction.
Consequently a level structure over a degree-five extension of $F$
has residue degree five. Its trace $D$ is a closed divisor.

Let $C_F$ be generated by the normalized global cusp-orbit divisors,
with rational base cusp $c_0$. There are $24$ orbit divisors,
of degrees $1^{16},2^8$. At $7$ the proved finite groups are
\[
\begin{aligned}
 J(\mathbb F_7)&\simeq[6,12,24,24,48,48],\\
 C_F&\simeq[6,12,12,12,24,24].
\end{aligned}
\]
In particular $\#C_F=5971968$.
For $A_{13}=T_{13}-1-13\langle13\rangle$, split-prime
Eichler--Shimura gives $A_{13}J(F)=0$.
The exact kernel in $J(\mathbb F_7)$ has quotient $C_2^2$
over the cusp image. Odd torsion injectivity gives
$2J(F)\subseteq C_F$; full cuspidal generation is not assumed.

The noncuspidal counts in degrees one through five at $7$ are
$(0,0,96,468,3264)$. Entirely cuspidal reduction is impossible.
We include all
\[
 96\left(\binom{17}{2}+8\right)+468\cdot16+3264=24576
\]
noncuspidal degree-five patterns.
The Hecke condition leaves $3552$ patterns, giving $3360$
distinct possible doubled global cusp classes. Thus
\[
 2D\sim L_v:=10c_0+\sum_{j=1}^{24}v_j(C_j-\deg(C_j)c_0)
\]
for an explicitly listed signed cusp vector $v$.
Eight completed native certificates recompute all $3828$ support
point classes, including the exact Hecke map on a generating set.
Independent integer-lattice replay checks every pattern and every lift.

We next reduce at $13$ and test the complete systems $|\overline L_v|$.
For speed, use the degree-six coordinate map $u:X\to\mathbb P^1$.
For a section $s$, the finite part of the pushforward of
$L_v+\operatorname{div}(s)$ is the divisor of a polynomial of
degree at most ten, obtained from the norm of $s$.
If $L_v+\operatorname{div}(s)$ is even, this polynomial is a
constant times a square. We reject nonsquares first and compute
full support for the remaining sections. Entirely cuspidal halves
are excluded by the same odd-prime argument.
The test examines $126480$ projective sections and leaves $192$
classes. Every rejected system is exhausted. All polynomial-square
and projective-coverage tests are replayed independently.
Each surviving system has dimension two at $13$, hence also
in characteristic zero by Riemann--Roch and specialization.

Ninety-six classes are
\[
 L_{8,i}+2c,\qquad 1\le i\le6,
\]
where the $c$ are the sixteen $F$-rational cusps and the $L_{8,i}$
are six effective cuspidal degree-eight divisors with
$\ell(L_{8,i})=2$. These are degree-independent pencil computations;
their equations and completed characteristic-zero certificates
are preserved with this proof.
The lattice of $141$ integral modular-unit divisors has rank $23$
and index $5971968$ in the degree-zero cusp lattice.
It verifies every displayed class identity exactly.
Since $\ell(L_{8,i}+2c)=2$, the complete system has $2c$ as a
fixed part. It cannot contain twice a closed quintic divisor.

The remaining $96$ classes form one orbit under the $96$ level maps
\[
 (P_6,Q_{12})\longmapsto(aP_6+2bQ_{12},cP_6+dQ_{12}),
\]
where $a,b,c$ are taken modulo $6$, $d$ modulo $12$,
$a,d$ are odd, and $ad-2bc\equiv1\pmod6$, modulo simultaneous sign.
Every image class is checked by fresh proved Jacobian arithmetic.
The equality of global classes follows by torsion injectivity at $7$.
It remains to consider one representative $L$.

Modular-unit lattice reduction gives an effective cuspidal
representative of degree ten. Its complete pencil has basis
$1,f$, with pole divisor $L$. Put
\[
 H(u,a)=h_L(u)\operatorname{Norm}_{F(X)/F(u)}(f-a),
\]
where $\operatorname{div}(h_L)=u_*L-10\infty$.
This is a polynomial of degree ten in $u$. An even fiber requires
its monic normalization to be a square.
Formal square completion, followed by the gcd of the remaining
coefficient conditions, gives
\[
 a-\frac{180+15z}{2}.
\]
The exceptional zeros of the leading coefficient are
\[
 \frac{181+13z}{2},\quad 89+8z,\quad
 \frac{179+17z}{2},\quad\frac{177+21z}{2}.
\]
At each of these four values the specialized polynomial is not
a constant times a square. The fiber at infinity is cuspidal.
Thus there is only one possible noncuspidal even fiber.

At $a=(180+15z)/2$ we obtain $2P$, where $P$ has degree five,
$u(P)=w$, and
\[
 v(P)=(4z+1)w^4+(4-z)w^3+(10-2z)w^2+(2-13z)w+2z-5.
\]
Here $w$ satisfies the quintic of Proposition~\ref{prop:decic612}.
The characteristic polynomial defining $H$, the square-completion
gcd, and all exceptional specializations are reconstructed in Sage.
Substitution into the curve equation and its implicit derivative
independently verifies the double zero at $P$.
Since the pencil has degree ten, the fiber is exactly $2P$.
This gives one closed quintic point in the representative class,
and the level automorphisms give all the others.
The full-torsion assertion follows from Proposition~\ref{prop:decic612}.
\end{proof}
